% !TeX root = ../../Book3.tex
% 纲要标题：### 23.2 平方零提升与形式性质
% 纲要位置：计划纲要.md，第 2710 行。

\section{平方零提升与形式性质}
\label{b3:sec:2302}

把误差放入平方零理想后，多项式方程的修正成为一个线性问题。提升的存在与唯一性分别约束关系模和微分模；有限展示另外保证这些条件能由有限组方程描述。

% 纲要内容（仅作写作计划，尚非教材正文）：
%
% * 平方零扩张的非空提升集合由导子加法群自由传递地作用；存在性与差的参数化分别判断
% * 形式光滑的关系模左逆判据；微分模自由仍可不形式光滑的正特征反例
% * 形式光滑、形式非分歧和形式平展的复合、基变换与目标局部化
% * 代数有限展示的作用及其与模有限展示的区别
%
% ---
%

\subsection{平方零扩张上的提升问题}
\label{b3:subsec:2302:01}

\chapref{b3:ch:14}已经把导子解释为平方零方向的变化。这里进一步追问：一个满足方程的近似点能否修正成精确点？令 \(A\to B\) 为交换含幺环同态，\(C\) 为交换 \(A\) 代数，\(J\subseteq C\) 为满足 \(J^2=0\) 的理想。给定 \(A\) 代数同态 \(u:B\to C/J\)，要求寻找 \(v:B\to C\)，使它与商映射的复合为 \(u\)。这就是本节的提升问题。

\ProseParagraphBreak
假如 \(B=A[X_1,\ldots,X_n]/(f_1,\ldots,f_s)\)，先任取 \(u(\bar X_i)\) 在 \(C\) 中的提升 \(c_i\)。这些元素未必满足关系；只知道每个误差 \(f_j(c)\) 属于 \(J\)。再以 \(h_i\in J\) 修正时，平方零条件使 Taylor 公式精确地变成
\[
f_j(c+h)=f_j(c)+\sum_i\frac{\partial f_j}{\partial X_i}(c)h_i.
\]
高次修正项全部为零。因此本来非线性的方程，变成一个关于误差的线性方程组。

\begin{lemma}\label{b3:lem:square-zero-lifts-difference}
设 \(A\to B\) 为交换含幺环同态，\(C\) 为交换 \(A\) 代数，\(J\subseteq C\) 为平方零理想，\(u:B\to C/J\) 为 \(A\) 代数同态，且至少存在一个提升 \(v:B\to C\)。由 \(u\) 使 \(J\) 成为 \(B\) 模。则全部提升恰为
\[
v+D,\qquad D\in\operatorname{Der}_A(B,J)
\cong\operatorname{Hom}_B(\Omega_{B/A},J).
\]
特别地，两个提升的差由微分模控制，而提升是否存在还需要另作判断。
导子的加法群在这个非空提升集合上自由且传递地作用；选定 \(v\) 才给出了参数化所用的原点。
\end{lemma}
\begin{proof}
\(C/J\) 对 \(J\) 的作用不依赖代表元，因为 \(J^2=0\)。若 \(v'\) 也是提升，则 \(D=v'-v\) 取值于 \(J\)，在 \(A\) 上为零；展开 \(v'(bb')\) 并删去两个 \(J\) 元素之积，得到
\[
D(bb')=u(b)D(b')+u(b')D(b).
\]
故 \(D\) 为导子。反之，同一个等式证明 \(v+D\) 保持乘法，也保持加法、单位与 \(A\) 的像，所以给出提升。最后一个同构是微分模的泛性质。
\end{proof}

存在一个提升以后，每个导子给出唯一的修正，两个提升之间也有唯一的差；在尚无提升时，这个导子模不保证方程有解。\(B=A[X]\) 时，任取一个变量像都能提升，而全部提升之间相差一个 \(J\) 中的元素。若 \(B=A[T]/(f)\) 且 \(f'(t)\) 可逆，则误差由唯一的修正 \(-f(c)/f'(c)\) 消去。前者有自由的变化方向，后者连这种方向也没有。

\subsection{形式光滑、形式非分歧与形式平展}
\label{b3:subsec:2302:02}

按照\cref{b3:def:formal-infinitesimal-properties}，对每个平方零提升问题都存在提升、至多有一个提升、恰有一个提升，分别称为形式光滑、形式非分歧、形式平展。本节沿用这些定义，不附加 Noether 条件。一般幂零理想 \(J^N=0\) 的问题可以沿 \(C/J,C/J^2,\ldots,C/J^N\) 逐层处理，因为每一步的核平方为零。

\begin{proposition}\label{b3:prop:formal-properties-stability}
交换含幺环同态的形式光滑性、形式非分歧性及形式平展性，分别在复合、任意基变换和目标环局部化下保持。
\end{proposition}
\begin{proof}
对复合 \(A\to B\to D\)，先将 \(B\to C/J\) 提升为 \(B\to C\)，从而使 \(C\) 成为 \(B\) 代数，再提升 \(D\) 的映射，便得存在性。若两个复合提升存在，形式非分歧性先使它们在 \(B\) 上相同，再使它们在 \(D\) 上相同，便得唯一性。

基变换 \(B'=A'\otimes_AB\) 的 \(A'\) 代数同态等同于保持既定 \(A'\) 作用的 \(B\) 的 \(A\) 代数同态。提升后由张量积泛性质延拓，存在性和唯一性都保留。最后，若 \(B\to S^{-1}B\)，一个在 \(C/J\) 中成为单位的元素，其任意提升在 \(C\) 中也为单位：若 \(cd=1+j\)、\(j\in J\)，则 \((1+j)^{-1}=1-j\)。因此对 \(B\) 的提升自动延拓到 \(S^{-1}B\)，且这种延拓唯一。
\end{proof}

\cref{b3:thm:formally-unramified-differentials}已证明，形式非分歧恰好等价于 \(\Omega_{B/A}=0\)。形式光滑需要处理的却是关系误差。对于 \(B=A[X_1,\ldots,X_n]/I\)，\(I/I^2\) 记录一阶关系，\(B^n\) 记录变量的修正方向，微分映射 \(\delta:I/I^2\to B^n\) 记录这些方向在各关系中的系数。若它有左逆 \(\tau\)，每个误差同态 \(\lambda:I/I^2\to J\) 就能扩为 \(\lambda\tau:B^n\to J\)，从而用变量修正消去误差。只研究余核 \(\Omega_{B/A}\) 是否投射，还不能保证这种扩张。

\begin{theorem}[形式光滑的关系模判据]\label{b3:thm:formal-smooth-conormal-splitting}\label{b3:thm:free-differentials-not-smooth}
设 \(A\) 为交换含幺环，\(P=A[X_1,\ldots,X_n]\)，\(I\subseteq P\) 为理想，\(B=P/I\)。令
\[
\delta:I/I^2\longrightarrow B^n,
\qquad \bar f\longmapsto
\left(\overline{\frac{\partial f}{\partial X_1}},\ldots,
\overline{\frac{\partial f}{\partial X_n}}\right).
\]
则 \(B\) 在 \(A\) 上形式光滑，当且仅当 \(\delta\) 有 \(B\) 线性的左逆。此时关系正合列
\[
0\longrightarrow I/I^2\xrightarrow{\delta}B^n
\longrightarrow\Omega_{B/A}\longrightarrow0
\]
分裂，特别地 \(\Omega_{B/A}\) 为有限生成投射 \(B\) 模。

这个后果不能反推形式光滑性。具体地，若 \(k\) 为特征二域，\(B_0=k[X]/(X^2)\)，则 \(\Omega_{B_0/k}=B_0\,\mathrm{d}\bar X\) 为秩一自由模，但 \(B_0\) 在 \(k\) 上不形式光滑，因而不光滑。
\end{theorem}
\begin{proof}
若形式光滑，将恒等映射 \(B\to B\) 提升为 \(s:B\to P/I^2\)。设 \(q:P\to P/I^2\)、\(\pi:P\to B\) 为自然映射，则
\(D=q-s\pi\) 取值于 \(I/I^2\)，并且由平方零条件是 \(A\) 导子。它在 \(f\in I\) 上取值 \(\bar f\)。令 \(\tau:B^n\to I/I^2\) 把第 \(j\) 个基向量送到 \(D(X_j)\)，则导子的多项式求导公式给出 \(\tau\delta=1\)。

反之，设有这样的 \(\tau\)。对任意 \(u:B\to C/J\)、\(J^2=0\)，任取变量像的提升，得到 \(w:P\to C\)。由 \(w(I)\subseteq J\)，限制映射给出 \(B\) 线性映射
\(\lambda:I/I^2\to J\)：乘以 \(P\) 中元素时，其作用只取决于在 \(B\) 中的像；而 \(I^2\) 映为零。设 \(h_j=(\lambda\tau)(e_j)\)，将第 \(j\) 个变量像改为 \(w(X_j)-h_j\)。对每个 \(f\in I\)，平方零 Taylor 公式使新像等于
\[
\lambda(\bar f)-\lambda\tau\delta(\bar f)=0.
\]
于是修正后的映射经过 \(B\)，构成所需提升。正合列的右半部分来自微分模的关系展示，左逆则保证左端单射及分裂。

对于最后的反例，取 \(P_0=k[X]\)、\(I_0=(X^2)\)。乘以 \(X^2\) 给出 \(B_0\cong I_0/I_0^2\)，所以关系模非零且自由秩一。但特征二使 \(\mathrm{d}(X^2)=0\)，故它到 \(B_0\,\mathrm{d}X\) 的映射为零，没有左逆；其余核 \(\Omega_{B_0/k}\) 仍是自由模。也可直接看提升失败：置 \(C=k[T]/(T^3)\)、\(J=(T^2)C\)，则 \(J^2=0\)、\(C/J\cong B_0\)。把 \(\bar X\) 送到 \(\bar T\) 的恒等识别若有提升，就须把 \(\bar X\) 送到 \(\bar T+h\)、\(h\in J\)；其平方却恒为非零的 \(\bar T^2\)，不能满足关系。
\end{proof}

这一证明的核心是由提升 \(B\to P/I^2\) 得到关系模的分裂；可对照松村英之\cite[Theorem 25.2, pp. 194--195]{Matsumura1989CommutativeRingTheory}。\(\Omega_{B/A}\) 投射只保证右端的满射 \(B^n\to\Omega_{B/A}\) 分裂，关系模判据还要求左边的 \(\delta\) 单射并有左逆；特征二的反例正是在这个左端失败。同一证明适用于 \(P\) 的局部化：变量像以外的要求只是将指定元素变成单位，而平方零提升自动保持这一点。

\subsection{有限展示条件的作用}
\label{b3:subsec:2302:03}

本卷约定：有限展示且形式光滑的映射称为光滑映射；有限型且形式非分歧的映射称为非分歧映射。平展映射目前定义为有限展示、平坦且 \(\Omega_{B/A}=0\)；下一节将证明它等价于有限展示且形式平展。这里说的是交换环同态 \(A\to B\) 的有限展示，即只要求有限个代数生成元和有限条关系；它与 \(B\) 作为 \(A\) 模有限展示是不同的条件。

\begin{proposition}\label{b3:prop:finite-presentation-filtered-maps}
设 \(A\) 为交换含幺环，\(B\) 为有限展示交换 \(A\) 代数，\((C_i)\) 为有向系统的交换 \(A\) 代数，\(C=\varinjlim_i C_i\)。则自然映射
\[
\varinjlim_i\operatorname{Hom}_{A\text{-alg}}(B,C_i)
\longrightarrow\operatorname{Hom}_{A\text{-alg}}(B,C)
\]
为双射。
\end{proposition}
\begin{proof}
写 \(B=A[X_1,\ldots,X_n]/(f_1,\ldots,f_s)\)。一个到 \(C\) 的映射只需指定有限个变量像，故它们在某个共同的 \(C_i\) 中已有代表。每个 \(f_j\) 在极限中为零，因而在某个较后的阶段为零。关系也只有有限条，再取一个共同阶段便得到到该阶段的映射，证明满射性。若两个阶段的映射在极限中相同，则有限个生成元的像分别在更后阶段相同；再次取共同阶段，两映射就相同，证明单射性。
\end{proof}

这个命题中，有限个代数生成元使变量像能够落在同一个阶段，有限条关系使它们确实在较后的同一阶段定义一个同态，而再比较有限个生成元的像便控制了两个同态的相等。下一节的局部标准形还会利用有限生成的关系理想，把一点处的关系等式扩大到一个主开邻域；标准模型的有限多个系数则允许把计算放到有限型整数代数上。因此有限展示同时给出有限阶段、邻域和系数环中的有限数据，形式提升性质本身不提供这些有限性。

\begin{example}\label{b3:exm:formal-without-finite-presentation}
设 \(A\) 为交换含幺环。无限变量多项式环 \(A[X_1,X_2,\ldots]\) 在 \(A\) 上形式光滑，因为可以逐个选取变量像的提升。但当 \(A\ne0\) 时，它不是有限型 \(A\) 代数：任意有限组多项式只涉及有限个变量，不能生成另外的变量。

\ProseParagraphBreak
另一个例子是 \(\overline{\mathbb Q}/\mathbb Q\)。每个有限子扩张均可分，可取本原元并用\cref{b3:prop:standard-etale-lift}得到形式平展性。对整个代数闭包的任一平方零问题，逐个有限子扩张都有唯一提升；唯一性使这些提升在交叠处一致，从而合成整个代数闭包的唯一提升。因此此扩张形式平展，却不是有限展示：有限个代数生成元只能生成有限次扩张。有限展示不能从“每次小范围提升都唯一”中省略。
\end{example}
